% ═══════════════════════════════════════════════════════════════════════════════
\chapter{Conclusion \& Future Work}
% ═══════════════════════════════════════════════════════════════════════════════

% intro comments, trim/ replace
This technical report details the design, implementation, and evaluation of a \gls{qkd} Simulator developed using Python, Qiskit, and MATLAB. By employing a modular \gls{oop} architecture, the project successfully modelled three \gls{dv}-\gls{qkd} protocols under varying channel conditions to evaluate the resilience of \gls{bb84}, \gls{b92}, and E91. The subsequent experiments translated abstract quantum mechanics into quantifiable engineering metrics, explicitly evaluating protocol viability for real-world deployment. This chapter evaluates the project's achievements against its initial objectives, synthesises key findings, and outlines the broader societal context and directions for future work.

\section{Achievement Against Objectives}
\label{sec:achievements}

The primary objective of developing a configurable, software-based \gls{qkd} testing framework was successfully achieved. The implementation of a YAML-based configuration system decoupled simulation parameters from the execution logic, enabling scalable batch testing across multiple execution modes. The core protocols were accurately modelled alongside a comprehensive post-processing pipeline capable of error estimation and Devetak-Winter key rate calculations. 


%  Maybe dont include a Gantt chart ask Martin
% replace gantt chart line with this: Time management strategies, facilitated by the modular \gls{oop} design, were critical in scoping the development.
Crucially, the project was completed in accordance with a structured software development methodology. Time management tools, such as Microsoft Planner, Toggl time-tracking and a Gantt Tracker, were critical in scoping the development of system features.


\section{Key Findings}
\label{sec:key}

The multivariable analysis produced several critical findings regarding protocol performance under interference:

\begin{itemize}
    \item \textbf{Architecture Validation:} Under ideal channel conditions, the simulation converged to the theoretical Mutual Information baseline of $1.0\ \text{bit/sifted qubit}$ with a $0.0\%$ \gls{qber}, validating the core cryptographic logic.
    \item \textbf{BB84 Resilience vs. \gls{b92} Fragility:} The four-state structure of \gls{bb84} demonstrated significant resilience to depolarising noise, maintaining a \gls{skr} of $\sim24 \text{ bits/100 qubits}$ at $5\%$ noise. Contrastingly, \gls{b92}'s two-state structure resulted in a severe sifting discard rate, causing its \gls{skr} to collapse to $\sim3 \text{ bits/100 qubits}$ under identical conditions.
    \item \textbf{Entanglement Degradation:} The simulation successfully mapped the threshold of \gls{e91} entanglement collapse. At a Noise Strength of $P \approx 0.165$, the S-parameter decayed from the quantum Tsirelson bound of $2\sqrt{2}$ to below the classical limit of $2$, resulting in key distribution failure.
    \item \textbf{Multivariable Deployment Viability:} The 3D surface analyses proved that while \gls{bb84} maintains a functional "secure operational volume" under simultaneous noise and interception, \gls{b92} rapidly breaches its $6.5\%$ \gls{qber} threshold. This suggests that the deployment of \gls{b92} is largely impractical for most environments.
\end{itemize}


\section{Evaluation of Methodology \& Limitations}
\label{sec:eval}


Whilst the simulator provides a rigorous framework for theoretical \gls{qkd} analysis, the reliance on Qiskit and NumPy abstracts certain physical imperfections such as detector inefficiencies, and fibre optic channel attenuation. The simulation also assumes trusted sources and detectors, artificially inflating key generation rates when compared to physical systems.


To mitigate the statistical uncertainty in quantum measurement simulations, the Central Limit Theorem and Monte Carlo approach were utilised to extract reliable metrics. As the channel interference models depict an approximation of physical systems, the derived key rates should be interpreted as theoretical upper bounds of operation.


\section{Broader Context, Ethics \& EDI}
\label{sec:context}

Developing accessible \gls{qkd} simulation tools holds significant importance within the broader context of network security. As classical cryptography becomes increasingly vulnerable to quantum threats, bridging the worlds of quantum information theory and practical network engineering is critical. Ethically, the transition to quantum-secure networks must prioritise data integrity and secure communication infrastructure for all sectors, not just wealthy actors.



From an \gls{edi} perspective, choosing to build the simulator using open-source libraries ensures the toolset remains accessible to researchers and students in all environments. By maintaining clear documentation and modular code architecture, the project lowers the barrier to entry for all demographics interested in quantum engineering.


\section{Future Work}
\label{sec:future}

Although the current version of the platform provides a thorough baseline for protocol analysis, several directions for future development exist:

\begin{itemize}
    % Point 1: \gls{e91} eve
    \item \textbf{Implementation of Eve for E91:} Unifying the benchmark by developing an optimal attack strategy against \gls{e91} would complete the comprehensive analysis of the three protocols.

     % Point 2: Expanding the protocol library
    \item \textbf{Implementation of Advanced Protocols:} Integrating alternative \gls{qkd} protocols described by Kumar \textit{et al.} \cite{Kumar2025} to expand the suite to modern protocols.

    % Point 3: Enhancing the noise modelling
    \item \textbf{Hardware Noise Integration:} Coupling the simulator with IBM Quantum's real hardware noise models to evolve the channel interference module from theoretical approximations to device-specific error characteristics.

    % Point 4: Scaling the simulation
    \item \textbf{Expanded Network Topology:} Advancing the architecture from simple \gls{p2p} links to multi-node quantum network topologies to enable the simulation of quantum repeaters and trusted-node routing.

\end{itemize}